% ==========================================================================
% GENERATED by scripts/split_paper.py. Do not edit this file.
%
% Source: imported_papers/grokking/iclr2027_conference_v2.tex
% Part 4: Whether a rule groks is decided at the minima of the loss
%
% Edits belong upstream, in the published paper. To change what the
% chapter includes, change imports/grokking.toml and re-run
% the splitter.
% ==========================================================================
\section{Whether a rule groks is decided at the minima of the loss}
\label{sec:direction}

\begin{table}[t]
\caption{The rules on both axes: backpropagation and the credit knob, our per unit rescaling of its credit magnitudes (Section~\ref{sec:select}), beside five of the twelve published rules (Table~\ref{tab:zoo} holds the rest). The Normalizes column gives which parameters share one normalizer in the rule's update, $\downarrow$ and $\uparrow$ the concentrated and the distributed end. Circuit PR and credit PR are the participation measures of Section~\ref{sec:select}, at $p=31$ over three seeds, standard deviations on circuit PR, credit PR a seed mean (Appendix~\ref{app:baselines}). A dash means not applicable: DFC and the feedback alignment rules do not rescale the gradient, a rule that builds no circuit has no circuit PR, and K-FAC here applies its own raw step (Appendix~\ref{app:consumer}).}
\label{tab:twoaxis}
\begin{center}
\scriptsize
\renewcommand{\arraystretch}{0.88}
\begin{tabular}{lccccl}
\toprule
Rule & Normalizes & Credit PR & Circuit PR & Test & Outcome \\
\midrule
LARS \citep{you2017lars} & per layer & $0.53$ & $0.63\,(0.04)$ $\downarrow$ & $0.970$ & groks \\
Backpropagation      & none           & $0.59$ & $0.60\,(0.02)$ $\downarrow$ & $0.996$ & groks \\
DFC                  & --- & $0.94$ & $0.92\,(0.00)$ $\uparrow$ & $1.000$ & groks \\
Knob ($\beta{=}0.25$)& per unit       & $0.96$ & $0.94\,(0.01)$ $\uparrow$ & $0.996$ & groks \\
Knob ($\beta{=}1$)   & per unit       & $1.00$ & $0.99\,(0.00)$ $\uparrow$ & $0.975$ & groks \\
\midrule
Sign-SGD             & per parameter  & $1.00$ & --- & $0.000$ & memorizes \\
K-FAC                & per unit       & $0.98$ & --- & $0.486$ & memorizes \\
DFA                  & ---            & $0.98$ & --- & $0.000$ & memorizes \\
FA (layerwise)       & ---            & $0.96$ & --- & $0.000$ & memorizes \\
\bottomrule
\end{tabular}
\end{center}
\end{table}

Changing the learning rule does not always preserve grokking. Backpropagation, DFC, the credit knob and all eight published optimizers grok, reaching test accuracy above $0.97$ (Table~\ref{tab:twoaxis}, with the optimizers in Table~\ref{tab:zoo}), while four rules never generalize at all and stay at chance on held out data however long they train. What separates the two groups is where training comes to rest. Once training accuracy reaches $1.0$, the parameters lie on $\mathcal{M} = \{\theta : L_{\mathrm{train}}(\theta) \approx 0\}$, the set of parameters that fit the training data, and stay inside it. We call it the \emph{interpolating set}. There the training loss is constant, so only the weight decay still pulls, and it pulls toward the solutions in $\mathcal{M}$ with the smallest weights. Those are the solutions that generalize, since memorizing needs larger weights than the periodic solution does \citep{liu2022omnigrok}.

Whether a rule comes to rest there is decided by the part of its update that is not the gradient. Split a rule's update into the gradient of the regularized loss and a remainder orthogonal to it, which we call the rule's \emph{transverse component} $\vb_t$. Training halts only where the full update vanishes, so a rule whose transverse component is nonzero at the minima of the loss cannot halt at one. It halts away from them, and the offset is largest along $\mathcal{M}$, where the loss is flattest and the circuit is decided, so training ends off the grokked circuit while training accuracy stays perfect (Proposition~\ref{prop:persist}, with the size of the offset in Appendix~\ref{app:theory}). Both parts shrink as training converges, so the size of $\vb_t$ on its own says little and we compare it against the gradient: the component \emph{survives} when the ratio $\|\vb_t\|/\|\vg_t\|$ does not fall as the loss does. This is the criterion we test: \emph{if a rule's transverse component vanishes as the loss falls, the network groks.}

For the rules that grok, the reason can be read from the update (Appendix~\ref{app:mech}). Backpropagation, LARS and the knob only rescale the gradient, so any orthogonal part they carry shrinks as the gradient does. DFC demonstrates that a rule can satisfy the criterion without staying close to the gradient: its update sits far from the gradient throughout training, and its transverse component still falls from $0.410$ of the gradient norm to $0.038$, because the controller signal driving it is proportional to the remaining output error, and training closes that error. A rule can therefore estimate the gradient badly and still grok, as long as the error fades where the loss does.

Each rule that fails to generalize holds a component that does not fade. Feedback alignment routes the error through a fixed random matrix, so its component keeps one direction and is still about a quarter of the gradient norm at the end. Sign-SGD's update has norm $\eta\sqrt{d}$ whatever the gradient, so at a fixed learning rate it cannot shrink relative to a vanishing gradient, and it still memorizes at every learning rate we swept, down to $10^{-5}$ (Appendix~\ref{app:named}). K-FAC fails for a different reason, and it shows that the credit assignment algorithm does not decide the outcome on its own: its own raw step memorizes, while the same credit handed to Adam groks (Appendix~\ref{app:consumer}). Whether a circuit forms is a property of the algorithm and the optimizer together, not of the credit alone.

The criterion holds without exception across these rules and three constructed controls, which also rule out three alternative explanations (Appendix~\ref{app:mech}). \emph{Persistence of direction} fails: DFC's transverse direction is almost as persistent as a fixed-bias control, autocorrelation $0.97$ against $0.99$, yet DFC groks and the control memorizes, and scaling that same control by the current loss so that it fades restores grokking everywhere it had destroyed it. \emph{Distance from the gradient} fails: at matched cosine, fresh noise leaves grokking intact down to $0.2$ while the fixed bias kills it by $0.7$. \emph{Credit magnitude} fails: the four rules that memorize spread credit \emph{more} evenly than backpropagation, and swapping feedback alignment's random matrix for the transpose of the forward weights recovers grokking, which validates our implementation. What selects a circuit is therefore not what decides whether one forms at all.
